\documentclass[11pt,a4paper]{article}
\usepackage[margin=1in]{geometry}
\usepackage{amsmath,amssymb}
\usepackage{parskip}
\usepackage{hyperref}

\title{\textbf{Pre-Experiment Overview:\\ Franck-Hertz Experiment}}
\author{IBRAHIM H.I. ABUSHAWISH \\ Istanbul University, Department of Physics}
\date{}

\begin{document}
\maketitle

\section*{Purpose of the Experiment}
According to the lab manual, the primary goal of this experiment is to prove the existence of discrete, quantized atomic energy levels by studying the inelastic collisions between free electrons and mercury atoms in a vapor-filled Franck-Hertz tube. By recording the anode current, $I_A$, as a function of accelerating voltage, $U_1$, we measure the spacing between successive current minima and calculate the first excitation energy of the mercury atom.

\section*{Importance and Significance}
Performed by James Franck and Gustav Hertz in 1914, this Nobel Prize-winning experiment provided the first direct, non-optical proof of the Bohr atomic model. Prior to this, evidence for discrete energy levels came exclusively from spectroscopic light absorption and emission. Franck and Hertz demonstrated that an atom cannot absorb arbitrary amounts of kinetic energy from colliding electrons; rather, energy transfer can only occur in discrete quantum packets matching the atom's internal excitation thresholds.

In the experiment, a sealed Franck-Hertz tube containing mercury droplet vapor is heated in an oven to approximately $160$--$180^\circ\text{C}$ to maintain an optimal collision density. Thermionically emitted electrons from a heated cathode are accelerated toward a positive perforated grid anode by voltage $U_1$, after which they must overcome a small retarding braking voltage, $U_2$, to reach the collector electrode. As $U_1$ increases, current rises until electrons acquire exactly $4.9\text{ eV}$ of kinetic energy---the threshold needed to excite a mercury valence electron from the $6^1S_0$ ground state to the $6^3P_1$ first excited state via inelastic collision. Having lost their kinetic energy, the electrons cannot overcome the counter-potential $U_2$, causing a sharp drop in collector current. As $U_1$ is increased further, this excitation process repeats periodically, creating a striking oscillatory curve with current drops spaced precisely $4.9\text{ V}$ apart. Measuring these equidistant intervals proves that atomic energy levels are quantized and allows us to calculate the wavelength of the corresponding UV photon emitted when the excited mercury atom relaxes.

\end{document}
